\section{Caso práctico: control de robots y multiagente}

\subsection{Elección entre PPO/SAC en robótica}

\begin{figure}[H]
\centering
\includegraphics[width=0.9\textwidth]{slide-18.jpg}
\caption{Robotics benchmark que muestra la diapositiva 18 de la clase: PPO se usa para un control on-policy estable, SAC destaca en espacios de acción continuos de alta dimensión.}
\end{figure}

En las tareas reales de control de robots, el entrenamiento suele enfrentar restricciones de muestreo y espacios de acción de alta dimensión. La arquitectura de trust-region más clipping que ofrece PPO facilita limitar el tamaño de paso en el equipo real, mientras que el buffer off-policy entropy-optimal de SAC permite un ajuste más rápido.

\begin{knowledgebox}{Recomendaciones para la programación en producción}
En tareas de robot arm / locomotion, una estrategia común es: usar pruebas de PPO con horizon corto para validar rápidamente la estructura de reward; migrar la configuración más prometedora a SAC, junto con soft target y automatic entropy tuning. Los resultados de PPO son más fáciles de interpretar y SAC ahorra más muestras.
\end{knowledgebox}

\subsection{Critic compartido en el entrenamiento multiagente}

Cuando varios agent comparten un entorno, un common critic o un centralized value es clave para mejorar la estabilidad. La práctica habitual es entrenar un joint critic $Q(\boldsymbol{s}, \boldsymbol{a})$ conjunto y luego actualizar el actor por partes:

$$\hat{A}_i = Q(\boldsymbol{s}, \boldsymbol{a}) - V(\boldsymbol{s})$$

El agent solo observa su propia action, pero el critic ve el estado global, lo que le permite juzgar con más precisión el valor de las decisiones locales.

\begin{warningbox}{Riesgos del critic compartido}
Si el critic depende en exceso de la global information, o no cuenta con una value decomposition adecuada, el gradient de un agent individual puede quedar enmascarado por el noise global. Las soluciones incluyen information bottleneck, agent-specific mask y periodic critic reset.
\end{warningbox}

\subsection{Resumen de la sección}

PPO y SAC desempeñan papeles distintos en el entrenamiento de agentes: PPO funciona como política segura, SAC gestiona el muestreo de alta frecuencia; en escenarios multiagente se necesita un critic compartido y prestar atención a la fuga de información.

